\documentclass[a4paper,11pt]{article}

% Lingua
\usepackage[T1]{fontenc}
\usepackage[italian]{babel}
\usepackage{microtype}
\usepackage[normalem]{ulem}
\usepackage{csquotes}
\usepackage{epigraph}
\setlength{\epigraphwidth}{0.5\textwidth}

% Font
\usepackage{libertinus}
% \usepackage{lmodern}

% Layout
\usepackage{geometry}
\geometry{
  a4paper,
  top=25mm,
  bottom=22mm,
  left=30mm,
  right=30mm
}

\setlength\parindent{0pt} % no spazio quando si va a capo

% Matematica
\usepackage{mathtools}
\usepackage{amsthm}
\usepackage[makeroom]{cancel}
\usepackage{amssymb}
\usepackage{faktor} % per il simbolo del quoziente

% Liste e figures
\usepackage{enumitem}
\usepackage{graphicx}
\usepackage{placeins}

\graphicspath{ {./images/} }

% Bibliografia
\usepackage[
  backend=biber,
  style=alphabetic
]{biblatex}

\addbibresource{references.bib}

% Hyperlinks e numerazione automatica
\usepackage[hidelinks]{hyperref}
\setcounter{tocdepth}{3}

% Colori per i commenti e box colorati
\usepackage{xcolor}
\usepackage[most]{tcolorbox}
\colorlet{myblue}{blue!70!black}
\colorlet{myred}{red!70!black}

\newtcolorbox{mystarthighlight}{
    enhanced, frame hidden,
    colback=blue!3, colupper=myblue,
    left=2pt, right=0pt,
    borderline west={1.5pt}{0pt}{myblue},
}

\newtcolorbox{myendhighlight}{
    enhanced, frame hidden,
    left=8pt, right=0pt,
    borderline west={1pt}{0pt}{black, dashed},
}

% PARAMETRI TCOLORBOX

    % enhanced,                   % abilita le funzionalità grafiche avanzate
    % breakable,                  % permette al box di andare su più pagine

    % COLORI
    % colback=red!3,              % colore dello sfondo
    % colframe=red!75!black,    % colore del bordo normale
    % colupper=black,           % colore del testo interno
    % coltitle=black,           % colore del titolo

    % BORDO
    % frame hidden,               % nasconde il bordo standard
    % borderline={0.8pt}{0pt}{red!75!black,dashed}, % spessore,distanza,colore
    % borderline east = " " "
    % borderline west = " " "
    % borderline nord = " " "
    % borderline sud = " " "
    % boxrule=0.8pt,            % spessore del bordo standard

    % ANGOLI
    % sharp corners,            % angoli completamente quadrati
    % rounded corners,          % angoli arrotondati
    % arc=2mm,                  % raggio degli angoli arrotondati

    % SPAZIATURA INTERNA
    % boxsep=5pt,               % spazio interno generale
    % left=8pt,                 % spazio interno a sinistra
    % right=8pt,                % spazio interno a destra
    % top=6pt,                  % spazio interno sopra
    % bottom=6pt,               % spazio interno sotto

    % SPAZIO ESTERNO
    % before skip=10pt,         % spazio prima del box
    % after skip=10pt,          % spazio dopo il box

    % TITOLO
    % title=Titolo,             % titolo del box
    % fonttitle=\bfseries,      % stile del titolo
    % coltitle=black,           % colore del titolo

    % DIMENSIONI
    % width=\textwidth,         % larghezza del box
    % height=4cm,               % altezza del box

    % ALLINEAMENTO DEL CONTENUTO
    % halign=left,              % allineamento orizzontale (left / center / right)
    % valign=center,            % allineamento verticale (top / center / bottom)

    % EFFETTI GRAFICI
    % drop shadow,              % aggiunge un'ombra al box


% Alcune macro
\newcommand{\Sym}{\text{Sym}}
\newcommand{\gen}[1]{\langle #1 \rangle}
\NewDocumentCommand{\ngen}{m g}{
  \langle\!\langle #1 \rangle\!\rangle_{\IfNoValueTF{#2}{}{#2}}
}
\newcommand{\pres}[2]{\langle #1\mid #2 \rangle}
\newcommand{\mylink}[1]{\href{#1}{[\textbf{link}]}}


%  Numerazione teoremi
\theoremstyle{plain}
\newtheorem{teo}{Teorema}[section]
\newtheorem{prop}[teo]{Proposizione}
\newtheorem{corollario}[teo]{Corollario}
\newtheorem{lemma}[teo]{Lemma}

\theoremstyle{definition}
\newtheorem{defin}[teo]{Definizione}
\newtheorem{esempio}[teo]{Esempio}
\newtheorem{esercizio}[teo]{Esercizio}

\newtheorem{dimostrazione}{Dimostrazione}
\newtheorem{oss}[teo]{Osservazione}


% Personalizzazione documento







\usepackage{float}
\NewDocumentCommand{\TabellaCorso}{mmmmmom}{

\begin{table}[H]
    \begin{mystarthighlight}
    \begin{tabular}{ll}
        \textbf{Nome Corso} & #1\\
        \textbf{Tipologia - Anno} & \IfBlankTF{#1}{--}{#2}\\
        \IfBlankTF{#3}{}{
            \textbf{Programma} & \mylink{#3}\\
        }
        \IfBlankTF{#4}{}{
            \textbf{Prove d'esame} & #4\\
        }
        \IfBlankTF{#5}{}{
            \textbf{Dispense Ufficiali} & \mylink{#5} \IfNoValueTF{#6}{}{(a.a. #6)}\\
        }
        \IfBlankTF{#7}{}{
            \textbf{Eserciziario Ufficiale} & \mylink{#7}\\
        }
    \end{tabular}
    \end{mystarthighlight}
\end{table}

}


\usepackage{array}
\usepackage{tabularx}
\usepackage{xparse}

\newcommand{\corso}[1]{\subsection{#1}}

\NewDocumentEnvironment{risorse}{+b}
{%
    \noindent
    \setlength{\tabcolsep}{2pt}%
    \begin{myendhighlight}
        \begin{tabularx}{\textwidth}{
            @{}
            >{\raggedright\arraybackslash}p{0.10\textwidth}
            >{\raggedright\arraybackslash}p{0.20\textwidth}
            >{\raggedright\arraybackslash}p{0.15\textwidth}
            >{\raggedright\arraybackslash}p{0.15\textwidth}
            >{\raggedright\arraybackslash}X
            @{}
        }
        #1
        \end{tabularx}%
    \end{myendhighlight}
}
{}

% \risorsa{autore}{formato}{anno}{link}{nota}
\NewDocumentCommand{\risorsa}{mmmm O{}}{%
    \mylink{#4} & #1 & #2 & #3 & #5 \\
}


\title{Organizzazione dei corsi e del materiale UniPi}
\author{Andrea Arielli}
\date{2026-2027}


\begin{document}

\maketitle
\newpage
\tableofcontents
\newpage


\section{Premessa}
Questo ha lo scopo di organizzare il materiale dei corsi, triennali e magistrali, del dipartimento di matematica dell'Università di Pisa, mettendo tutte le informazioni utili allo studio in un unico posto.
Le informazioni possono non essere complete e/o corrette. Non garantisco che i link continuino a rimanere validi nel tempo, io farò quel che posso per mantenere il file aggiornato. 

Saranno elencate diverse branche (Algebra, Probabilità, Analisi, Geometria, ecc.) e per ognuna di queste verranno elencati i corsi. Per ogni corso indicherò dove trovare il materiale o talvolta inserirò un link. A volte inserirò solo il link alla pagina web del corso (o del professore che lo tiene nell'anno corrente) dove uno può trovare da solo diverse informazioni. L'a.a. corrente è 2026-2027, tuttavia sono presenti link a materiali relativi a corsi degli anni precedenti.

Alcune info generiche su dove trovare materiale in autonomia: per quanto riguarda le note/dispense ufficiali delle lezioni, se previste, le si trovano sulla pagina elearning del corso, o sulla pagina team del corso oppure sul sito web personale del docente (spesso è \url{https://people.dm.unipi.it/} + cognome del professore).
Per quanto riguarda gli appunti, manoscritti o in \LaTeX, basta cercare ``\texttt{poisson phc appunti [nome\_corso]}'' su google.


\newpage

\section{Corsi di Algebra}

\corso{Aritmetica 1}
\TabellaCorso{Aritmetica 1}{Triennale - 1}{https://elearning.dm.unipi.it/pluginfile.php/40957/mod_resource/content/1/programma26-27.pdf}{Su elearning} {https://elearning.dm.unipi.it/pluginfile.php/36230/mod_resource/content/20/dispenseAritmetica2024.pdf}[2025-2026] {https://elearning.dm.unipi.it/pluginfile.php/38623/mod_resource/content/1/Soluzioni_Aritmetica.pdf}

\begin{risorse}
    % \risorsa{Ilaria del Corso}{Ufficiali}{\LaTeX}{2025-2026}{}
    \risorsa{Simone Saccani}{Manoscritto}{2022-2023}{https://poisson.phc.dm.unipi.it/~saccani/materiale/Aritmetica.pdf}
    \risorsa{Diego Monaco}{\LaTeX}{2021-2022}{https://github.com/diego-unipi/Aritmetica/blob/main/Appunti\%20Aritmetica.pdf}
    \risorsa{Luca Bruni}{\LaTeX}{2015-2016}{https://poisson.phc.dm.unipi.it/~lbruni/Appunti/Aritmetica.pdf}
\end{risorse}
% \subsubsection*{Info Ulteriori}


\subsection{Algebra 1}
\TabellaCorso{Algebra 1}{Triennale - 2}{https://elearning.dm.unipi.it/pluginfile.php/37426/mod_resource/content/2/progrA1_2526.pdf}{Su elearning}{}{}

\begin{risorse}
    \risorsa{Simone Saccani}{Manoscritto}{2023-2024}{https://poisson.phc.dm.unipi.it/~saccani/materiale/Algebra\%20I.pdf}
    \risorsa{Daria Pasqualetti}{\LaTeX}{2023-2025}{https://poisson.phc.dm.unipi.it/~pereira/AlgebraI/Algebra1.pdf}
    \risorsa{Diego Monaco}{\LaTeX}{2022-2023}{https://github.com/diego-unipi/Algebra-1/blob/main/Appunti\%20Algebra\%201.pdf}
    \risorsa{Luca Bruni}{\LaTeX}{2016-2017}{https://poisson.phc.dm.unipi.it/~lbruni/Appunti/Algebra1.pdf}
\end{risorse}
\paragraph{Info ulteriori.}Pagina del professore, viene aggiornata e contiene diverse info, \mylink{https://people.dm.unipi.it/sbarra/pages/Didattica/aal1-2026/info_e_orario-aal1-2026.html}.
\subsection{Algebra 2}
\TabellaCorso{Algebra 2}{Magistrale/Triennale - 1/3}{}{}{https://people.dm.unipi.it/sbarra/pdfs/Didattica/al2-2021/A2-0.9.9.2-20210625.pdf}{}


\begin{risorse}
    \risorsa{Luca Bruni}{\LaTeX}{2017-2018}{https://poisson.phc.dm.unipi.it/~lbruni/Appunti/Algebra2.pdf}
    \risorsa{Carlo Siricana}{\LaTeX}{?-?}{https://poisson.phc.dm.unipi.it/~sircana/alg2.pdf}
\end{risorse}

\subsection{Istituzioni di Algebra}

\TabellaCorso{Istituzioni di Algebra}{Magistrale - 1}{https://unipi.coursecatalogue.cineca.it/corsi/2026/11514/insegnamenti/2026/53888_707507_77460/2026/53888?schemaid=9893}{}{https://poisson.phc.dm.unipi.it/~lmigliorini/primo_magistrale/istalg/istalg_lombardo.pdf}[2024-2025]{}

\paragraph{Info ulteriori.} Corso tenuto ogni anno da Daniele Lombardo, \mylink{https://people.dm.unipi.it/lombardo/} alla sua pagina web. Sulla pagina elearning del corso si trovano diversi testi di esercizi, su quella degli anni precedenti vi sono anche note manoscritte delle varie lezioni e testi di prove di esame.


\newpage
\section{Corsi di Probabilità}

\subsection{Elementi di probabilità}
\begin{risorse}
    \risorsa{Luca Bruni}{\LaTeX}{2017-2018}{https://poisson.phc.dm.unipi.it/~lbruni/Appunti/EPS.pdf}
\end{risorse}

\subsection{Probabilità}
\TabellaCorso{Probabilità}{Triennale - 3}{}{Su elearning}{https://elearning.dm.unipi.it/course/section.php?id=6749}[2026-2027]{}
\begin{risorse}
    \risorsa{Luca Bruni}{\LaTeX}{2018-2019}{https://poisson.phc.dm.unipi.it/~lbruni/Appunti/Probabilita.pdf}
\end{risorse}

\newpage
\section{Corsi di Analisi}

\subsection{Analisi Matematica 2}
\TabellaCorso{Analisi Matematica 2}{Triennale - 2}{}{\mylink{https://web.dm.unipi.it/paolini/didattica/2025/analisi2/}}{https://codeberg.org/paolini/AnalisiDue/releases}[2026-2027]{}
\begin{risorse}
    \risorsa{Simone Saccani}{\LaTeX}{2023-2024}{https://poisson.phc.dm.unipi.it/~saccani/materiale/Analisi_II.pdf}
    \risorsa{Francesco Sorce}{\LaTeX}{2022-2023}{https://www.overleaf.com/read/zgkydyyvcfyb\#bb45fd}
\end{risorse}

\paragraph{Info ulteriori.} La pagina web del professore contiene una pagina \mylink{https://web.dm.unipi.it/paolini/didattica/2026/analisi2/} dedicata al corso di Analisi 2 per l'a.a. 2026-2027.

\subsection{Analisi Matematica 3}
\TabellaCorso{Analisi Matematica 3}{Triennale - 3}{https://unipi.coursecatalogue.cineca.it/corsi/2024/10299/insegnamenti/2026/52834_694521_71577/2008/52834?schemaid=9144}{}{https://drive.google.com/file/d/14uRA4O-PRKiaLsngD-lfb_4tVZvc1dyv/view}[2026-2027]{}
\begin{risorse}
    \risorsa{Andrea Rocca}{\LaTeX}{2024-2025}{https://poisson.phc.dm.unipi.it/~rocca/materiali/analisi3.pdf}
    \risorsa{Simone Saccani}{manoscritto}{2024-2025}{https://poisson.phc.dm.unipi.it/~saccani/materiale/Analisi\%20III.pdf}
    \risorsa{Emanuele  Angile}{\LaTeX}{2018-2019}{https://poisson.phc.dm.unipi.it/~angile/analisi3.}
\end{risorse}

\paragraph{Info ulteriori.} Altre info si trovano sulla pagina teams del corso.

\subsection{Istituzioni di Analisi}
\begin{risorse}
    \risorsa{Francesco Sorce}{\LaTeX}{2024-2025}{https://github.com/fsorce/Istituzioni-di-analisi/blob/main/main.pdf}
\end{risorse}

\newpage
\section{Corsi di Sistemi Dinamici}

\subsection{Sistemi Dinamici}
\TabellaCorso{Sistemi Dinamici}{Magistrale/Triennale - 1/3}{}{Sul sito del professore Bonanno.}{https://pagine.dm.unipi.it/bonanno/Compiti/dyn-syst-book-public.pdf}[2026-2027]{}
\begin{risorse}
    \risorsa{Francesco Sorce}{\LaTeX}{2023-2024}{https://github.com/fsorce/Sistemi-Dinamici_2324/blob/main/main.pdf}
\end{risorse}
\paragraph{Info ulteriori.} Corso tenuto dai professori Claudio Bonanno \mylink{https://pagine.dm.unipi.it/bonanno/} e -, qui c'è il \mylink{https://pagine.dm.unipi.it/bonanno/sistdin-2627.html} alla pagina web del corso

\subsection{Teoria Ergodica}

\subsection{Meccanica Superiore}

\newpage
\section{Link utili e da aggiungere}
\begin{enumerate}
    \item [\mylink{https://poisson.phc.dm.unipi.it/~tateo/appunti/appunti.html}] contiene diversi appunti per triennale e magistrale formato manoscritto, 2013-2018 circa, ci sono diversi esercizi e tracce di esami vecchi
    \item[\mylink{https://web.evanchen.cc/napkin.html}] diverse dispense formato \LaTeX, lingua inglese, credo non unipi ma vastissimo
    \item[\mylink{https://mattiapuddu25.github.io/LectureNotes.html}] diverse dispense formato \LaTeX
    \item [\mylink{https://poisson.phc.dm.unipi.it/~frascella/appunti.html}] diverse dispense formato manoscritto, Frascella, copre molti corsi
    \item [\mylink{https://poisson.phc.dm.unipi.it/~glaudo/FilesFolder/GammaFunction.pdf}] Note sulla funzione $\Gamma$, tutte le dimostrazioni sono scritte in modo abbastanza elementare da studenti di Analisi 1 di Unipi. Introducono e raccolgono tantissimi risultati su questa funzione.
    \item [\mylink{https://poisson.phc.dm.unipi.it/~glaudo/FilesFolder/HausdorffDistance.pdf}] Note sulla distanza di Haussdorff. Sono risultati esposti da tre normalisti in un mini-seminario che tenuto come complemento al corso ordinario di matematica del primo anno. Nel documento si dimostra, partendo dalle definizioni, le prime proprietà della distanza di Hausdorff sullo spazio dei chiusi e limitati di un metrico.

    \item [\mylink{https://github.com/dario2994/CoolMathProblems/tree/master}] Raccolta di esercizi molto particolari di uno studente della normale della triennale su github.
\end{enumerate}



\newpage

% \printbibliography

\end{document}